Nous l'avons déjà évoqué, comprendre la notion de symétrie se fait avec la théorie des groupes, leurs actions et leurs représentations. Jusqu'à présent, nous nous sommes essentiellement intéressés aux groupes finis pour lesquels nous avons développé un ensemble de méthodes algébriques et combinatoires. \\

Cependant, de nombreuses situations conduisent à étudier des symétries continues. Les groupes qui les décrivent, tels que le groupe des rotations du plan $SO(2)$ ou de l'espace $SO(3)$, possèdent en plus de leur structure de groupe une structure différentielle lisse compatible. De tels objets sont appelés groupes de Lie.  \\

L'un des apports fondamentaux de cette théorie est d'exploiter la structure différentiable afin d'en extraire une description infinitésimale. En se plaçant au voisinage de l'élément neutre, on associe à tout groupe de Lie une algèbre de Lie, définie comme son espace tangent en l'élément neutre, muni d'une structure algèbrique supplémentaire : le crochet de Lie. \\

Cette correspondance entre un groupe de Lie et son algèbre de Lie est fondamentale : la structure locale du groupe de Lie est entièrement déterminée par son algèbre de Lie. Ainsi cela nous fournira une méthodologie afin d'étudier les groupes de Lie : on étudie d'abord son algèbre de Lie puis on se rammène au groupe de Lie via cette correspondance. \\

Par ailleurs, l’étude intrinsèque des algèbres de Lie révèle une structure remarquablement rigide. En particulier, dans le cas complexe, les algèbres de Lie semi-simples admettent une classification complète, reposant sur des données combinatoires finies : les diagrammes de Dynkin.